\subsection{A noncyclic even disk remains nonexact after its full constant quotient}
\label{arch136:noncyclic-constant-quotient}

\paragraph{Archive and exact scope.}
This tests the nonlinear quotient used in cote136, physical
p.~117/footer116/scan117 and physical p.~121/footer120/scan121.
For a commutative coefficient ring \(A\) and an \(A\)-module \(M\),
write
\[
 D_A(M)=\Gamma_A(M[0])\otimes_A\Lambda_A(M[1]),
 \qquad dm=e_m,
 \qquad
 Q_A(M)=D_A(M)/J_A(M),                                  \tag{NC.1}
\]
where \(J_A(M)\) is the full differential PD ideal generated by
all closed zero-forms in the augmentation kernel of
\(D_A(M)\to A\).  Parameter arity counts both an even generator
and its odd differential as one.  Coefficients have arity zero.
This is precisely the full extra-constant quotient, not merely
an ideal generated by selected examples.

The elementary calculation below is independent of any external
Poincare lemma.  The actual-tail application uses the explicit
resolution and strict-retract proof of
ARCH-136-ACTUAL-TAIL-SYZYGY, which were separately read in full.
No conclusion about the subsequent complete
AQ-136-RELATIVE-H1 repair is asserted.

\paragraph{An arity observation.}
Let \(A\) be a \(\mathbb Z_{(p)}\)-algebra.  On an even disk
the derivation \(h(e_m)=m\), \(h(m)=0\), of form degree minus one,
satisfies
\[
                       dh+hd=r\operatorname{id}
                   \quad\hbox{in parameter arity }r.   \tag{NC.2}
\]
This holds on arbitrary modules: the derivation respects every
linear relation and the special divided-power rule; the identity
is checked on generators and then on their products and divided
powers.  Thus no extra closed zero-form exists in positive arity
less than \(p\).  Every generator of \(J_A(M)\) has arity at least
\(p\).  Its form-one part in arity \(p\) is zero: a product with a
form-one element has arity at least \(p+1\), all higher divided
powers increase arity, and the generators themselves are closed.
Consequently
\[
                   H^1(D_A(M))_{(p)}
                    \xrightarrow{\ \sim\ }
                   H^1(Q_A(M))_{(p)}.                 \tag{NC.3}
\]
Indeed the degree-one terms and differentials into degree two
are unchanged in this arity, and the only degree-zero elements
being quotiented out are already closed.

\paragraph{The finite one-relation detector.}
Let \(k\) be a field, let
\[
 A=k[X,Y]/(X,Y)^2,
 \qquad M=(Au\oplus Av)/(Xu+Yv).                       \tag{NC.4}
\]
In parameter arity \(n\geq2\), there is an exact calculation
\[
                    H^1(D_A(M))_{(n)}
                         \simeq (k/nk)^{n-2}.           \tag{NC.5}
\]
In particular, if \(k=\mathbb F_p\), \(p\geq3\), then
\[
                         F=X\gamma_{p-1}(u)e_v          \tag{NC.6}
\]
is nonzero in this cohomology, and by (NC.3) remains nonzero
in the full quotient \(Q_A(M)\).

\begin{proof}
Put \(r=Xu+Yv\), so \(dr=Xe_u+Ye_v\).
The differential PD ideal of the linear relation is the ordinary
differential ideal generated by \(r,dr\).  All higher divided
powers of \(r\), and of its multiples, vanish already in the
free algebra because \((X,Y)^2=0\).  Splitting according to
coefficient degree zero or one, the degree-zero coefficient part
is the free even disk on two variables and is exact in every
positive arity.  All cohomology in question therefore comes
from the coefficient span \(kX\oplus kY\).

Write \(D_\mathrm{free}\) for the free two-variable disk over
\(k\), in the fixed positive arity \(n\).  Its sequence
\[
 0\longrightarrow D_\mathrm{free}^0
 \xrightarrow d D_\mathrm{free}^1
 \xrightarrow d D_\mathrm{free}^2\longrightarrow0
                                                               \tag{NC.7}
\]
is exact by the integral even-disk contraction.  In the quotient
(NC.4), the entire coefficient-degree-one part in form two is
zero: \(dr\,e_v=Xe_ue_v\) and
\(dr\,e_u=-Ye_ue_v\), and multiplying these identities covers
its divided-power basis.  Hence every coefficient-degree-one
one-form is closed.

By (NC.7), first quotienting such one-forms by the differentials
of all free zero-forms identifies their space with
\((kX\oplus kY)\otimes D_\mathrm{free}^2\).
Let
\[
 A_j=\gamma_j(u)\gamma_{n-2-j}(v)e_ue_v,
                         \qquad 0\leq j\leq n-2.
\]
The remaining relations are the differentials of
\(dr\) times free zero-forms and of \(r\) times free one-forms.
The former give exactly
\[
 XA_0=0,\qquad YA_{n-2}=0,\qquad
                 XA_j=YA_{j-1}\quad(1\leq j\leq n-2).
                                                               \tag{NC.8}
\]
For example
\(d(dr\,\gamma_a(u)\gamma_b(v))
=YA_{a-1}-XA_a\), where \(a+b=n-1\), with the absent
endpoint terms omitted.  After (NC.8) the free remaining basis
is \(YA_0,\ldots,YA_{n-3}\).

For \(a+b=n-2\), differentiating
\(r\gamma_a(u)\gamma_b(v)e_u\) gives
\[
                   -(b+1)YA_a-(a+1)XA_{a+1},
\]
and differentiating the same expression with \(e_v\) gives
\[
                    (a+1)XA_a+(b+1)YA_{a-1}.
\]
Together with (NC.8), these are precisely the relations
\(nYA_j=0\), \(0\leq j\leq n-3\).  There are no additional
relations, because all degree-one generators of the original
ideal have just been listed.  This proves (NC.5).

Under the isomorphism induced by \(d\) in (NC.7), the class
(NC.6) maps to
\(XA_{p-2}=YA_{p-3}\), one of the surviving basis vectors.
This proves its nonvanishing.  The coefficient-degree splitting
also shows that it cannot acquire a boundary from the free
coefficient-degree-zero part.  Finally (NC.3) applies to the
full extra-constant quotient.
\end{proof}

\paragraph{ARCH-136-NONCYCLIC-CONSTANT-QUOTIENT.}
For every odd prime \(p\), put \(R=\mathbb Z_{(p)}\) and
\(S=\Gamma_R(T)\).  There is a weightwise \(R\)-free,
\(S\)-torsion-free module \(M\), and also a bounded lattice
\(M_0=P+p^{-1}Q\) with two successive saturated free-kernel
presentations, for which \(H^1(Q_S(M))\), respectively
\(H^1(Q_S(M_0))\), contains a nonzero class of exact order
\(p\).  A single weightwise finite example can have such classes
in unbounded weights.

\begin{proof}
Fix \(h=p^a\), and put
\[
 x=T^{[h]},\quad y=T^{[ph]},\quad
 x^p=p\kappa y,\quad\kappa\in R^\times,
 \qquad M=(x,y)\subset S.
\]
Use the generators \(u=y\), \(v=-x\), of weights \(ph,h\).
Their complete relation module is generated by
\[
              (x,y),\qquad (p,\kappa^{-1}x^{p-1}).       \tag{NC.9}
\]
Indeed \(a y-bx=0\) is equivalent to
\(a x^{p-1}=p\kappa b\).  In \(S/pS\), the coefficient
\(\bar x\) is one truncated-polynomial digit generator,
and its \((p-1)\)-st power has annihilator \((\bar x)\).
Thus \(a=pc+xd\) and
\(b=\kappa^{-1}x^{p-1}c+yd\), proving (NC.9).

The underlying ring map \(S\to A\) of (NC.4) retains the two
digit generators \(x,y\) as \(X,Y\), and kills all other
digits, after reduction modulo \(p\).  Since \(p\geq3\),
\(X^{p-1}=0\).  By the full presentation (NC.9) and right
exactness of base change,
\[
                 M\otimes_S A\simeq(Au\oplus Av)/(Xu+Yv).
                                                               \tag{NC.10}
\]
The ring map is not asserted to preserve coefficient divided
powers.  The underlying relative module divided-power/exterior
functor still gives a differential algebra map, preserving the
new-variable PD operations and the relative augmentation.
It preserves the PD operations on the whole relative augmentation
ideal, which contains the constant ideal under discussion.  It
therefore takes every extra closed zero-form and its full PD ideal
to the corresponding target ideal, and induces a map on (NC.1).

The element
\[
             f=x\gamma_{p-1}(u)e_v                     \tag{NC.11}
\]
is closed: the relation \(xe_u+ye_v=0\), wedged with
\(e_v\), gives \(xe_ue_v=0\).  Moreover
\[
                 d\bigl(x\gamma_{p-1}(u)v\bigr)=pf.    \tag{NC.12}
\]
To verify the coefficient, multiply \(xe_u+ye_v=0\) by
\(\gamma_{p-2}(u)v\), and compare it with
\((xu+yv)\gamma_{p-2}(u)e_v=0\).  This gives
\(x\gamma_{p-2}(u)v e_u=(p-1)f\), and hence (NC.12).
Its image in (NC.10) is the nonzero (NC.6), including after
the full quotient.  Thus (NC.11) has exact order \(p\) in
\(H^1(Q_S(M))\).  Its weight is
\((p^2-p+2)h\).  The ideal \(M\subset S\) is visibly
\(S\)-torsion-free and weightwise \(R\)-free.

For the claimed lattice, let
\(P=S f_u\oplus S f_v\twoheadrightarrow M\),
\(Q=\ker(P\to M)\), and \(M_0=P+p^{-1}Q\).
It contains
\[
 \widetilde u=f_u-
       \frac{p f_u+\kappa^{-1}x^{p-1}f_v}{p}
          =-\frac{x^{p-1}}{p\kappa}f_v,
 \qquad \widetilde v=f_v.                               \tag{NC.13}
\]
These satisfy \(x\widetilde u+y\widetilde v=0\) and map
to \(u,v\).  Thus the same closed expression, relation
(NC.12), and detector prove the result for \(M_0\).
The quotient \(P/Q=M\) is torsion-free over \(S\), so
\(Q\) is saturated.  In addition \(M\) is isomorphic,
by its first coordinate, to
\[
          \ker\bigl((p,x):S\oplus S\langle h\rangle\to S\bigr).
\]
This kernel has image \((p,x)\), a torsion-free ideal, as
its quotient in the displayed free module, so that embedding
is saturated too.  Explicitly the inverse first-coordinate map
sends \(m\in M\) to \((m,-pm/x)\); it is integral by
(NC.9).  Hence these are two actual successive saturated kernels.
Weightwise \(M_0\subset p^{-1}P\) is free over the PID \(R\).

Direct sums over arbitrarily large \(h=p^a\), with a common
positive shift in each block if desired, have finitely many
contributors in any weight.  The inclusion and projection of
each block induce a retract of disks and also of the full
quotients, because maps preserve the relative augmentation and
closed zero-forms.  Thus the classes survive in one module,
in unbounded weights.  A weight shift \(v\) adds \(pv\)
to the weight of (NC.11).
\end{proof}

\paragraph{ARCH-136-ACTUAL-TAIL-CONSTANT-QUOTIENT.}
Let \(p\geq3\), \(N=p^a\), and let
\(P_\bullet\to I_N\) be any graded free \(S\)-resolution
with generator weights at least \(N\).  For even \(D\geq4\),
put
\[
 P=P_{D-2},\qquad Q=\ker(P_{D-2}\to P_{D-3}),\qquad
                         M_0=P+p^{-1}Q.                 \tag{NC.14}
\]
Then \(H^1(Q_S(M_0))\) contains nonzero order-\(p\) classes
in unbounded weights.  Thus the full extra-constant quotient
alone does not repair the actual even disk used at the bottom
of the coherent resolution.  This theorem concerns the standalone
even disk on the actual bottom parameter module.  It does not assert
a nonzero class in the full coherent realization: its additional
form-one and higher generators may kill that class.

\begin{proof}
Use first the explicit minimal resolution of
ARCH-136-ACTUAL-TAIL-SYZYGY.  In its notation
\(x_i=\gamma_{p^i}(T^{[N]})\), \(h_i=Np^i\), and
\(x_i^p=p\kappa_i x_{i+1}\).  Its free complex resolving the
augmentation has
\(\partial e_i=x_i\),
\(\partial f_i=x_i^{p-1}e_i-p\kappa_i e_{i+1}\), with
homological degrees one and two.  Write
\(r=(D-2)/2\), \(x=x_i\), \(y=x_{i+1}\), and
\(\kappa=\kappa_i\).  In its degree \(D-1\) free module
\(P_{D-2}\), take
\[
             \alpha=\gamma_r(f_i)e_i,
             \beta=\gamma_r(f_i)e_{i+1}.
\]
Let \(L=\operatorname{im}(P_{D-2}\to P_{D-3})\), and put
\(u=\partial\beta\), \(v=-\partial\alpha\).  The explicit
odd-image calculation gives
\[
                   xu+yv=0,
            \qquad pu+\kappa^{-1}x^{p-1}v=0.             \tag{NC.15}
\]
There is a detector \(L\otimes_S A\to M_A\), with \(A,M_A\)
as in (NC.4), sending \(\partial\beta\) to \(u\),
\(-\partial\alpha\) to \(v\), and every other free
presentation generator to zero.  We verify all its relations,
which is stronger than the pure-PD detector used previously.

After retaining only \(x_i,x_{i+1}\) in the square-zero
coefficient ring, every coefficient of \(\partial f_j\)
vanishes: it is either divisible by \(p\) or a
\((p-1)\)-st power of a radical element.  Only the Koszul
terms \(\partial e_j=x_j\) remain.  A term of an incoming
differential can therefore project onto \(\alpha\) or
\(\beta\) only if its divided-power part is
\(\gamma_r(f_i)\).  The only source with two retained
exterior indices is
\(\gamma_r(f_i)e_ie_{i+1}\); its projection is
\(x\beta-y\alpha\).  The resulting target relation is
exactly \(Xu+Yv\).  Sources with any other removed exterior
index have zero coefficient; all divided-power derivative
terms were already zero.  This verifies every incoming
relation and proves the detector.

In the actual lattice (NC.14), use
\[
       \widetilde u=\frac{x^{p-1}}{p\kappa}\alpha,
                 \qquad \widetilde v=-\alpha.            \tag{NC.16}
\]
The difference of \(\widetilde u\) from \(\beta\) is
in \(p^{-1}Q\), by the second relation of (NC.15).
They map to \(u,v\) and satisfy
\(x\widetilde u+y\widetilde v=0\) literally.
Equations (NC.11)--(NC.12) and the new detector give the
nonzero class in \(H^1(Q_S(M_0))\).  Its weight is
\[
               \bigl(rp^2+p^2-p+2\bigr)h_i,             \tag{NC.17}
\]
which tends to infinity with \(i\).  A common parameter
weight shift adds \(pv\), as before.

Finally the explicit minimal resolution is a strict chain
retract of every allowed free resolution, as proved in
ARCH-136-ACTUAL-TAIL-SYZYGY by comparison maps and the finite-weight
triangular inverse.  The induced maps on bottom lattices are
still a retract.  Both the disk and its full extra-constant
quotient are functorial, so the same retract detects the
classes for any allowed resolution.  This proves the assertion.
\end{proof}

\paragraph{Exact remaining gate.}
The cyclic quotient theorem cannot extend to arbitrary noncyclic
modules, even within the actual coefficient-tail lattice class.
The first positive class above survives the full constant quotient
for an arity reason, not merely because selected constants were
missed.  Its degree is one.  AQ-136-RELATIVE-H1 is explicitly
designed to kill such classes, and this calculation makes no
claim that the class survives that further repair.  Nor does it
identify all new degree-two or degree-three classes after that
repair, or show that an admissible common source is impossible.

\paragraph{Verification status.}
The author has completed a full rereading of this saved module,
including the exact arity calculation, full constant ideal, integral
module presentation, bounded lattice, and every incoming relation
in the actual-tail detector.  The parent and a contributing peer
have each completed a full saved-proof reading and reported PASS.  A finite matrix probe suggested
the one-relation example;
the proof above uses no extrapolation from that probe.  The all-prime
calculation and actual-tail detector were developed with the
cote15 branch; the separate records distinguish discussions from
complete saved-proof reviews.
